
\documentclass{article}
\usepackage{colortbl}
\usepackage{makecell}
\usepackage{multirow}
\usepackage{supertabular}

\begin{document}

\newcounter{utterance}

\centering \large Interaction Transcript for game `cladder', experiment `full\_v1.5\_default', episode 9292 with qwen3.5{-}9b.
\vspace{24pt}

{ \footnotesize  \setcounter{utterance}{1}
\setlength{\tabcolsep}{0pt}
\begin{supertabular}{c@{$\;$}|p{.15\linewidth}@{}p{.15\linewidth}p{.15\linewidth}p{.15\linewidth}p{.15\linewidth}p{.15\linewidth}}
    \# & \multicolumn{2}{c}{Player} && \multicolumn{2}{c}{Game Master} \\
    \hline

    \theutterance \stepcounter{utterance}  
    & & & \multicolumn{4}{p{0.6\linewidth}}{
        \cellcolor[rgb]{0.9,0.9,0.9}{
            \makecell[{{p{\linewidth}}}]{
                \texttt{\tiny{[P1$\langle$GM]}}
                \texttt{You are an expert in causal inference. The following question is not a typical commonsense query, but rather a meticulously designed question created by a professor specializing in causal inference, intended to assess the students' mastery of the course content.} \\
\texttt{Imagine a self{-}contained, hypothetical world with only the following conditions, and without any unmentioned factors or causal relationships: Zory has a direct effect on zuph and glimx. Zuph has a direct effect on jyka. Jyka has a direct effect on glimx. Zory is unobserved. The overall probability of zuph is 36\%. The probability of not zuph and glimx is 30\%. The probability of zuph and glimx is 31\%. Is the chance of glimx smaller when observing zuph?} \\
\texttt{Start your answer with 'Yes' or 'No', followed by additional reasoning or evidence to support your explanation.} \\
            }
        }
    }
     \\ \\

    \theutterance \stepcounter{utterance}  
    & \multicolumn{4}{p{0.6\linewidth}}{
        \cellcolor[rgb]{0.561,0.737,0.561}{
            \makecell[{{p{\linewidth}}}]{
                \texttt{\tiny{[P1$\rangle$GM]}}
                \texttt{No} \\
\\ 
\texttt{**Reasoning:**} \\
\\ 
\texttt{To determine if the chance of `glimx` is smaller when observing `zuph`, we must calculate the conditional probability \$P(\textbackslash{}text\{glimx\} \textbackslash{}mid \textbackslash{}text\{zuph\})\$ and compare it to the marginal probability \$P(\textbackslash{}text\{glimx\})\$.} \\
\\ 
\texttt{1.  **Calculate the marginal probability of `glimx` (\$P(\textbackslash{}text\{glimx\})\$):**} \\
\texttt{    We are given the joint probabilities involving `zuph` and `glimx`:} \\
\texttt{    *   \$P(\textbackslash{}neg \textbackslash{}text\{zuph\} \textbackslash{}land \textbackslash{}neg \textbackslash{}text\{glimx\}) = 0.30\$} \\
\texttt{    *   \$P(\textbackslash{}text\{zuph\} \textbackslash{}land \textbackslash{}text\{glimx\}) = 0.31\$} \\
\texttt{    *   \$P(\textbackslash{}text\{zuph\}) = 0.36\$} \\
\\ 
\texttt{    First, find \$P(\textbackslash{}text\{zuph\} \textbackslash{}land \textbackslash{}neg \textbackslash{}text\{glimx\})\$:} \\
\texttt{    \$\$P(\textbackslash{}text\{zuph\} \textbackslash{}land \textbackslash{}neg \textbackslash{}text\{glimx\}) = P(\textbackslash{}text\{zuph\}) {-} P(\textbackslash{}text\{zuph\} \textbackslash{}land \textbackslash{}text\{glimx\}) = 0.36 {-} 0.31 = 0.05\$\$} \\
\\ 
\texttt{    Next, find \$P(\textbackslash{}neg \textbackslash{}text\{zuph\} \textbackslash{}land \textbackslash{}text\{glimx\})\$:} \\
\texttt{    Since the total probability of the universe is 1, and we know \$P(\textbackslash{}neg \textbackslash{}text\{zuph\} \textbackslash{}land \textbackslash{}neg \textbackslash{}text\{glimx\}) = 0.30\$, the remaining probability mass belongs to cases where `zuph` is false but `glimx` is true (or vice versa, but we just calculated the `zuph` true cases).} \\
\texttt{    Alternatively, sum the probabilities where `glimx` is true:} \\
\texttt{    \$\$P(\textbackslash{}text\{glimx\}) = P(\textbackslash{}text\{zuph\} \textbackslash{}land \textbackslash{}text\{glimx\}) + P(\textbackslash{}neg \textbackslash{}text\{zuph\} \textbackslash{}land \textbackslash{}text\{glimx\})\$\$} \\
\\ 
\texttt{    We know \$P(\textbackslash{}neg \textbackslash{}text\{zuph\} \textbackslash{}land \textbackslash{}neg \textbackslash{}text\{glimx\}) = 0.30\$.} \\
\texttt{    The sum of all four joint states must be 1:} \\
\texttt{    \$\$P(\textbackslash{}text\{zuph\} \textbackslash{}land \textbackslash{}text\{glimx\}) + P(\textbackslash{}text\{zuph\} \textbackslash{}land \textbackslash{}neg \textbackslash{}text\{glimx\}) + P(\textbackslash{}neg \textbackslash{}text\{zuph\} \textbackslash{}land \textbackslash{}text\{glimx\}) + P(\textbackslash{}neg \textbackslash{}text\{zuph\} \textbackslash{}land \textbackslash{}neg \textbackslash{}text\{glimx\}) = 1\$\$} \\
\texttt{    \$\$0.31 + 0.05 + P(\textbackslash{}neg \textbackslash{}text\{zuph\} \textbackslash{}land \textbackslash{}text\{glimx\}) + 0.30 = 1\$\$} \\
\texttt{    \$\$0.66 + P(\textbackslash{}neg \textbackslash{}text\{zuph\} \textbackslash{}land \textbackslash{}text\{glimx\}) = 1\$\$} \\
\texttt{    \$\$P(\textbackslash{}neg \textbackslash{}text\{zuph\} \textbackslash{}land \textbackslash{}text\{glimx\}) = 0.34\$\$} \\
\\ 
\texttt{    Now, calculate the marginal \$P(\textbackslash{}text\{glimx\})\$:} \\
\texttt{    \$\$P(\textbackslash{}text\{glimx\}) = 0.31 + 0.34 = 0.65\$\$} \\
\\ 
\texttt{2.  **Calculate the conditional probability \$P(\textbackslash{}text\{glimx\} \textbackslash{}mid \textbackslash{}text\{zuph\})\$:**} \\
\texttt{    \$\$P(\textbackslash{}text\{glimx\} \textbackslash{}mid \textbackslash{}text\{zuph\}) = \textbackslash{}frac\{P(\textbackslash{}text\{zuph\} \textbackslash{}land \textbackslash{}text\{glimx\})\}\{P(\textbackslash{}text\{zuph\})\}\$\$} \\
\texttt{    \$\$P(\textbackslash{}text\{glimx\} \textbackslash{}mid \textbackslash{}text\{zuph\}) = \textbackslash{}frac\{0.31\}\{0.36\} \textbackslash{}approx 0.861\$\$} \\
\\ 
\texttt{3.  **Comparison:**} \\
\texttt{    *   \$P(\textbackslash{}text\{glimx\}) = 0.65\$} \\
\texttt{    *   \$P(\textbackslash{}text\{glimx\} \textbackslash{}mid \textbackslash{}text\{zuph\}) \textbackslash{}approx 0.861\$} \\
\\ 
\texttt{    The probability of `glimx` increases from 0.65 to approximately 0.86 when `zuph` is observed. Therefore, the chance of `glimx` is **not** smaller; it is larger.} \\
\\ 
\texttt{    *(Note: The causal structure described—Zory \$\textbackslash{}to\$ zuph, Zory \$\textbackslash{}to\$ glimx, zuph \$\textbackslash{}to\$ jyka, jyka \$\textbackslash{}to\$ glimx—implies a confounding relationship between `zuph` and `glimx` due to the common cause `Zory`. Since `Zory` is unobserved, observing `zuph` provides information about `Zory`, which in turn updates the probability of `glimx`. The calculation confirms this positive association.)*} \\
            }
        }
    }
    & & \\ \\

    \theutterance \stepcounter{utterance}  
    & & & \multicolumn{2}{p{0.3\linewidth}}{
        \cellcolor[rgb]{0.9,0.9,0.9}{
            \makecell[{{p{\linewidth}}}]{
                \texttt{\tiny{[GM$|$GM]}}
                \texttt{No} \\
            }
        }
    }
    & & \\ \\

    \theutterance \stepcounter{utterance}  
    & & & \multicolumn{2}{p{0.3\linewidth}}{
        \cellcolor[rgb]{0.9,0.9,0.9}{
            \makecell[{{p{\linewidth}}}]{
                \texttt{\tiny{[GM$|$GM]}}
                \texttt{no} \\
            }
        }
    }
    & & \\ \\

    \theutterance \stepcounter{utterance}  
    & & & \multicolumn{2}{p{0.3\linewidth}}{
        \cellcolor[rgb]{0.9,0.9,0.9}{
            \makecell[{{p{\linewidth}}}]{
                \texttt{\tiny{[GM$|$GM]}}
                \texttt{game\_result = WIN} \\
            }
        }
    }
    & & \\ \\

\end{supertabular}
}

\end{document}
